%% ma: L10-B15-0091
%% lop: 10
%% chuong: 6
%% bai: 15
%% dang: 10.15.4
%% loai: TN4
%% muc: TH
%% dap_an: "A"
%% nguon: "Ảnh giáo viên gửi 25/11/2026 (lô 2): Bài tập trắc nghiệm – Toán 10 Bài 15 Hàm số (sự biến thiên, đồ thị), Phần I Câu 2"
%% kien_thuc: "Lớp 10 Bài 15 (sự đồng biến, nghịch biến của hàm số)"
%% trang_thai: chinh-thuc
\begin{ex}
Xét sự biến thiên của hàm số $f(x)=\dfrac3x$ trên khoảng $(0;+\infty)$. Khẳng định nào sau đây đúng?
\choice{\True Hàm số nghịch biến trên khoảng $(0;+\infty)$.}{Hàm số vừa đồng biến, vừa nghịch biến trên khoảng $(0;+\infty)$.}{Hàm số đồng biến trên khoảng $(0;+\infty)$.}{Hàm số không đồng biến, không nghịch biến trên khoảng $(0;+\infty)$.}
\loigiai{Với $0<x_1<x_2$: $f(x_1)-f(x_2)=\dfrac{3(x_2-x_1)}{x_1x_2}>0$. Vậy hàm số nghịch biến trên $(0;+\infty)$.}
\end{ex}
